
\fchapter{Pág 112}

\fsection[\textcolor{waveBlue2}{Ejercicio 9}]{\boldit{\textcolor{waveBlue2}{9.}}{
Elige un producto que consumas o utilices en tu vida diaria. Razona si cumple o no con los principios del ecodiseño. ¿Podrías sustituir este producto por una opción más sostenible?
}}

\begin{solution}
\textbf{Producto: Cinturón de levantamiento Strengthshop}

\begin{itemize}
    \item \textbf{Cumplimiento de principios de ecodiseño:}  
    El cinturón Strengthshop está fabricado con materiales duraderos (cuero sintético o cuero reforzado y hebillas metálicas), garantizando larga vida útil y fiabilidad durante entrenamientos de fuerza. Esto cumple parcialmente con los principios de ecodiseño relacionados con \textbf{durabilidad y resistencia}. Mejoras:
        \begin{itemize}
            \item La hebilla se puede sustituir en caso de romperse (\textasciitilde30 €), evitando reemplazar el cinturón completo. Esto aumenta la sostenibilidad al reducir residuos.
        \end{itemize}
    Aspectos que podrían mejorar:
        \begin{itemize}
            \item No incorpora materiales reciclados ni reciclables.
            \item No hay indicios de que la producción minimice el consumo de energía o residuos.
        \end{itemize}

    \item \textbf{Posible sustitución por opción más sostenible:}  
    Se podría optar por cinturones fabricados con materiales reciclados o de origen local, producidos con procesos de bajo impacto ambiental. Sin embargo, dado que la hebilla es reemplazable, mantener y reparar el cinturón Strengthshop actual ya mejora significativamente su sostenibilidad.
\end{itemize}
\end{solution}
